\documentclass[11pt,oneside]{amsart}
\usepackage[margin=1in]{geometry}
\usepackage{times}
\usepackage{graphicx}
\usepackage{url}
\parindent=0pt
\pagestyle{empty}

\newcommand{\header}[1]{\bigbreak\textbf{#1}}

\begin{document}

\begin{center}
  \bf
  \includegraphics[width=3in]{logo} \\
  Syllabus \\
  STAT 695 \\ 
 Competitve Predictive Modeling  \\
  Spring, 2016 \\
  (2 credits)                 
\end{center}

\header{Description}

Statistical machine learning has became the central tool for a large number of research and practical fields, such as business decision making, imaging processing, and detecting disease relevant factor, and particularly predictive modeling. A vast amount of statistical tools and models have been discussed in literature for predictive modeling from both theoretical and methodological perspectives. In this course, instead of focusing on the theoretical aspects, students will gain extensive practice in building and testing predictive models through directed participation in predictive modeling competitions, including Kaggle (\url{https://www.kaggle.com/}) and the Data Mining Cup (\url{http://www.data-mining-cup.de/en/}). Competitions will include predicting responses of mixed types. Practical utilization of statistical learning tools will be discussed along with the competition. In addition, students will gain extensive practice in using R
software for data wrangling and modeling. 


\header{Prerequisites}

Experience coding in R; at least one regression course such as STAT 511; permission of the instructors.


\header{Instructors}

Professor Brooke Anderson, \url{brooke.anderson@colostate.edu}

Office: 146 Environmental Health Building

Office hours: By appointment \medskip

Professor Wen Zhou, \url{riczw@colostate.edu}, \\
\url{http://www.stat.colostate.edu/~riczw/}

Office: 208 Statistics Building

Office hours: By appointment



\header{Meetings}

Tuesdays 4:00-4:45; Thursdays 3:50-4:05; Location: Stats 006



\header{Textbooks}

Required: 
\begin{itemize}
\item Max Kuhn and Kjell Johnson, \emph{Applied Predictive
  Modeling}, Springer (2013). (Available online through CSU library)
 
 \item Gareth James, Daniela Witten, Trevor Hastie, Robert
Tibshirani, \emph{An Introduction to Statistical Learning: With
  Applications in R}, Springer (2013). \url{http://www-bcf.usc.edu/~gareth/ISL/ISLR%20First%20Printing.pdf}
\end{itemize}

Recommended: 
\begin{itemize}
\item Alan J. Izenman, {\it Modern Multivariate Statistical Techniques: Regression, Classification, and Manifold Learning}, Springer (2008).
 
 \item Trevor Hastie, Robert Tibshirani, and Jerome Friedman, {\it The Elements of Statistical Learning: Data Mining, Inference, and Prediction, Second Edition}, Springer (2009) (The book is free online).

\item Christopher M. Bishop, {\it Pattern Recognition and Machine Learning}, (2006).
\item Ian Witten, Eibe Frank, and Mark Hall, {\it Data Mining: Practical Machine Learning Tools and Techniques, Third Edition}, (2011).
\end{itemize}





\header{Course Objectives}

The primary aim of this course is to use three in-depth, practical
examples of predictive-modeling challenges (two Kaggle competitions
and the Data Mining Cup) to allow students to gain extensive practical
experience in data processing, data matrix construction, variable selection, model fitting, model ensemble, and evaluating predictive
models. With that aim, this course will introduce students to a
variety of topics in machine learning to provide them with strategies
and approaches to tackle these specific applied examples. Some of the
challenges that will come up will be specific to these three problems,
but this experience of completing three predictive modeling
competitions will provide students with the conceptual understanding, programming tools, and
strategies to tackle their own predictive modeling challenges. This
course would also provide an excellent background for students aiming
to take more theoretical courses on statistical learning in the
future. 
\vspace{.2cm}

Specific objectives include: 

\begin{enumerate}
\item Students will work in multi-disciplinary teams to compete in the two Kaggle predictive modeling competitions.
\item Students will learn how to fit and evaluate a variety of
  predictive models, including: classification and regression trees,
  support vector machines, logistic and linear regression models,
  tree ensemble models, Naive Bayes, k-nearest neighbors, and neural networks.
\item Students will learn strategies in data wrangling and feature engineering to improve
  predictive models.
\item Students will learn to use resampling methods to assess the
  performance of predictive models.
\item Students will gain extensive additional experience working on
  complex modeling problems using R statistical software. 
  \item Students will be recommend to form a team to participate the annual Data Mining Cup held from April to May. The final result will be released on July 2016, and the top 10 teams will be invited to Berlin, Germany, for the final ceremony; more details see \url{http://www.data-mining-cup.de/en/}, \url{http://magazine.amstat.org/blog/2013/10/01/iowa-state-dmc/}.
\end{enumerate}


\header{Course Topics}

\begin{itemize}
\item Definition of machine learning.
\item Classification models: K-nearest neighbors, naive Bayes,
  logistic regression, classification trees, bagging, random forests,
  boosting, support vector machines.
\item Regression models: K-nearest neighbors, naive Bayes, linear
  regression, regression trees and tree ensembles, non-linear
  regression, support vector machines, ridge regression, lasso.
\item Model fitting and tuning.
\item Model evaluation, including with re-sampling techniques.
\item Data pre-processing, feature selection, measuring variable
  importance, and visualizing data.
\item Linear model selection and regularization and the challenges of
  high-dimensional data.
\item Neural networks and deep learning.
\item Unsupervised learning.
\end{itemize}


\header{Course Expectations \& Grading}

Grades will be based on attendance, participation, regular submission
of model results to each of the two Kaggle competitions, written reports on
each challenge (one per group), in-class presentations on final
models for each Kaggle challenge and the DMC, and student presentations on topics from
Kuhn and Johnson (2013).

\header{Assignments \& Readings}

\begin{itemize}
\item \textbf{Week 1:} What is machine learning? Classification models: K-nearest neighbors
  and naive Bayes.\\
\emph{Reading:} James et al.: Chs. 2, 4.1-2, 4.4. Kuhn and Johnson
Chs. 1, 2, 13.5-6.\\
\emph{Competition:} Kaggle: Surviving the Titanic

\item \textbf{Week 2:} Metrics of performance of classification
  models. Classification models: logistic regression models and
  classification trees.\\
\emph{Reading:} James et al.: Chs. 2.2, 4.3, 4.5, and 8.1. Kuhn and
Johnson: Chs. 11, 12.1-3, 14.1-2.\\
\emph{Tuesday assignment:} Fit some different Naive Bayes and k-Nearest Neighbors models. For a few of your models, try measuring some other metrics of performance, like sensitivity, specificity, AUC, Youden's J Index, positive predictive value, etc.\\ 
\emph{Thursday assignment:} Try fitting the following types of models: LDA (see week 1 reading), logistic regression, classification tree. Can you get any of your models to beat the Sex-only benchmark model? Can you get any to beat the to-date class best (Casey's 0.8)?\\
\emph{Competition:} Kaggle: Surviving the Titanic

\item \textbf{Week 3:} Using resampling to measure performance of
  classification models. Classification models: ensemble models
  (bagging, random forest, boosting).\\
\emph{Reading:} James et al.: Chs. 5, 8.2. Kuhn and Johnson: Ch. 14.3-7.\\
\emph{Tuesday assignment:} Try using cross-validation, either to tune certain paramters in a model (e.g., k in k-NN), or to compare the performance of different models\\ 
\emph{Thursday assignment:} Try fitting the following types of ensemble models: bagging, boosting, random forests. Can you get any to beat the to-date class best (~0.8)?\\
\emph{Competition:} Kaggle: Surviving the Titanic

\item \textbf{Week 4:} Classification models: support vector
  machines.\\
\emph{Reading:} James et al.: Ch. 9\\
\emph{Tuesday assignment:} Try to fit a SVM model to the Titanic data\\
\emph{Thursday assignment:} Present 5 slides on your work on the Titanic prediction as well as a short write-up describing your work on the challenge.\\
\emph{Competition:} Kaggle: Surviving the Titanic\\
\textbf{Graded products:} Students present final predictive model for Titanic
competition.

\item \textbf{Week 5:} Student presentations based on Kuhn and Johnson
  2013: Data pre-processing; Over-fitting and model tuning; Remedies
  for severe class imbalance; Feature selection\\
\emph{Reading:} Kuhn and Johnson: Chs. 3, 4, 16, 19\\
\emph{Tuesday:} Presentations from Casey, Veronica, and Wande\\
\emph{Thursday:} Presentations from Julia and Ahmed. Discussion of all presentations\\
\textbf{Graded products:} Students presentations on material from chapters in Kuhn and
Johnson 2013. Presentations are (all from Kuhn and Johnson): Ch. 3.1-3.3: Casey; Ch. 3.4-3.7: Veronica; Ch. 4.6-4.8: Wande; Ch. 16.1-16.7: Julia; 19.1-3: Ahmed\\

\item \textbf{Week 6:} Assessing performance of regression
  models. Regression models: Linear regression models. Over-fitting
  and model tuning re-visited.\\
\emph{Reading:} James et al.: Ch. 3. Kuhn and Johnson: Chs. 4
(re-visited), 5, 6.1-2.\\
\emph{Competition:} Kaggle: Current competition with continuous
outcome

\item \textbf{Week 7:} Linear model selection and
  regularization. Shrinkage methods and dimension reduction methods.\\
\emph{Reading:} James et al.: Ch. 6.\\
\emph{Competition:} Kaggle: Current competition with continuous
outcome

\item \textbf{Week 8:} Regression models: Non-linear regression
  models, regression trees. Measuring predictor importance.\\
\emph{Reading:} James et al.: Ch. 7. Kuhn and Johnson: Chs. 7 and 8.\\
\emph{Competition:} Kaggle: Current competition with continuous
outcome

\item \textbf{Week 9:} Regression models: Regression trees. Feature selection /
  engineering re-visisted.\\
\emph{Reading:} James et al.: Ch. 8 (re-visited). Kuhn and Johnson:
Chs. 3 and 19 (re-visited).\\
\emph{Competition:} Kaggle: Current competition with continuous
outcome

\item \textbf{Week 10:} Student team presentations on final model for current Kaggle
  competition.\\
\emph{Reading:} Background information on this year's Data Mining Cup\\
\emph{Competition:} Kaggle: Current competition with continuous
outcome\\
\textbf{Graded products:} Students present final predictive model for
current Kaggle competition.\\

\item \textbf{Week 11:} Factors that can affect model performance. Case study: Grant application models. \\
\emph{Reading:} Kuhn and Johnson: Chs. 15, 20. \\
\emph{Competition:} This year's Data Mining Cup competition.

\item \textbf{Week 12:} High-dimensional data (re-visited). Case study: Concrete Mixture Strength models.\\
\emph{Reading:} James et al. : Ch. 6 (re-visited). Kuhn and Johnson: Ch. 10.\\
\emph{Competition:} This year's Data Mining Cup competition.

\item \textbf{Week 13:} Neural networks. Deep learning methods. \\
\emph{Reading:} Kuhn and Johnson: Ch. 13.2. \\
\emph{Competition:} This year's Data Mining Cup competition.

\item \textbf{Week 14:} Visualization. Unsupervised learning.\\
\emph{Reading:} James et al.: Ch 10. \\
\emph{Competition:} This year's Data Mining Cup competition.

\item \textbf{Week 15:} Student team presentations on final model for
  this year's DMC  competition.\\
\emph{Competition:} This year's Data Mining Cup competition.\\
\textbf{Graded products:} Students present final predictive model for
DMC.

\end{itemize}



\end{document}
